\chapter{Fourier Intuition}
\label{ch:ch19-fourier-intuition}

\section{Periodic Cycles and Hidden Patterns}

In ancient astronomy, one of the most useful discoveries was the
\emph{Metonic cycle}: nineteen solar years are almost exactly equal to
$235$ lunar months. Neither cycle is defined in terms of the other; one
measures the motion of the Earth around the Sun, and the other measures
the phases of the Moon. Yet when the two cycles are observed together, a
larger repeating pattern emerges. This observation suggests a broader
question --- what happens when several periodic processes are combined?
--- and the answer leads to one of the most influential ideas in
mathematics, science, and engineering: complicated periodic behavior can
often be understood as a combination of simpler periodic motions.

\section{From Simple Oscillations to Complex Patterns}

Throughout this book we have encountered periodic functions such as
$\sin x$ and $\cos x$, functions that repeat indefinitely. Many
phenomena in nature are also periodic --- musical tones, ocean tides,
planetary motion, electrical current, and seasonal variation all recur
in cycles --- yet most real signals are not perfect sine waves. A violin
note is not a pure sinusoid, ocean tides are not perfectly smooth, and
human speech contains many interacting oscillations at once. This raises
a natural question: can complicated periodic patterns be built from
simple ones?

\section{Frequency Revisited}

In Chapter~\ref{ch:ch06-graphs}, we saw that $\sin(nx)$ oscillates more
rapidly than $\sin x$.

\begin{definition}
The \emph{frequency} of a periodic oscillation measures how rapidly it
repeats. For trigonometric functions, $\sin(nx)$ and $\cos(nx)$ have
frequency $n$ times that of $\sin x$ and $\cos x$.
\end{definition}

For example, $\sin(3x)$ completes three full oscillations in the same
interval where $\sin x$ completes one.

\section{Harmonics}

Frequencies that are whole-number multiples of a fundamental frequency
are called harmonics.

\begin{definition}
If $\sin x$ is taken as a fundamental oscillation, then $\sin(2x),
\sin(3x), \sin(4x), \dots$ and $\cos(2x), \cos(3x), \cos(4x), \dots$ are
called its \emph{harmonics}.
\end{definition}

Musical instruments rarely produce a single frequency; instead, they
produce combinations of harmonics, and the particular mixture of
harmonics present helps determine the distinctive sound, or
\emph{timbre}, of an instrument.

\section{Superposition}

The simplest way to create a more complicated oscillation is to add
simple oscillations together. Consider $f(x) = \sin x + \tfrac12\sin(2x)$:
the resulting graph is periodic, but it is no longer a pure sine wave.
Adding another harmonic produces even greater complexity, as in $f(x) =
\sin x + \tfrac12\sin(2x) + \tfrac14\sin(3x)$; each component remains
simple, but their combination produces a richer pattern. This principle
is called \emph{superposition}.

\begin{definition}
\emph{Superposition} is the formation of a new function by adding
simpler functions together.
\end{definition}

The central idea of Fourier analysis is that surprisingly complicated
periodic functions can often be represented through superposition of
harmonics.

\section{Complex Exponentials and Harmonic Sums}

Chapter~\ref{ch:ch16-complex-euler} introduced complex exponentials,
$e^{i\theta} = \cos\theta + i\sin\theta$, and
Chapter~\ref{ch:ch18-sequences-series} showed how finite sums of complex
powers may be analyzed algebraically. Consider a finite sum
\[
F(\theta) = \sum_{k=-n}^{n} c_k e^{ik\theta}.
\]
Each term $e^{ik\theta}$ represents a rotating unit vector whose angular
speed is proportional to $k$, so $F(\theta)$ is a finite superposition of
rotating vectors, in exactly the sense developed in the previous chapter.
Using Euler's formula, $e^{ik\theta} = \cos(k\theta) + i\sin(k\theta)$,
so
\[
F(\theta) = \sum_{k=-n}^{n} c_k\bigl(\cos(k\theta) + i\sin(k\theta)\bigr).
\]
When the coefficients satisfy $c_{-k} = \overline{c_k}$, the imaginary
contributions cancel in pairs, leaving a real-valued function.
Consequently, finite sums of complex exponentials correspond precisely
to finite sums of sine and cosine harmonics: the rotating vectors of
Chapter~\ref{ch:ch16-complex-euler} become the oscillating harmonics of
this chapter, related by nothing more than Euler's formula already in
hand.

\section{The Fourier Idea}

The central claim of Fourier analysis is remarkably simple: many
periodic functions can be represented as sums of harmonics,
\[
f(x) = a_0 + \sum_{n=1}^{\infty} \bigl(a_n\cos(nx) + b_n\sin(nx)\bigr),
\]
where the coefficients $a_n$ and $b_n$ determine how much of each
harmonic is present. This statement is one of the most important ideas
in applied mathematics. In a more advanced course, calculus provides
methods for computing the coefficients themselves; here the focus is on
the idea alone, that complicated periodic behavior may be assembled from
simple oscillatory components.

\section{Why Orthogonality Matters}

Chapter~\ref{ch:ch13-vectors-dot-products} studied orthogonality using
the dot product: two perpendicular vectors represent independent
directions. Fourier analysis relies on a closely related idea, in which
different harmonics behave like independent directions in a much larger
space, and the coefficient associated with a harmonic measures how
strongly that harmonic contributes to the overall signal. Just as a
vector may be decomposed into components along coordinate axes, a
periodic function may be decomposed into components along harmonic
directions. The analogy is not exact, but it captures the essential
geometric intuition: Fourier coefficients play a role similar to vector
projections.

\section{Visualizing Decomposition}

A useful analogy comes from optics. White light appears uniform, yet
when passed through a prism it separates into many component colors; the
light was not changed by the prism, which merely revealed structure
already present within it. Fourier analysis plays a similar role: a
complicated waveform may appear indivisible, yet Fourier decomposition
reveals the frequencies hidden within it, and musical tones, speech
signals, ocean waves, and electrical currents may all be viewed through
this same frequency-based perspective.

\begin{practicenow}
Identify the frequencies present in $f(x) = 2\sin x + \cos(3x)$. Which
frequency contributes the larger amplitude?
\end{practicenow}

\section{Applications}

Fourier analysis appears throughout science and engineering, wherever
oscillatory phenomena are present: in music and acoustics, where it
explains timbre; in audio compression, filtering, and signal processing
generally; in the study of heat flow and diffusion, the problem that
originally motivated it; in astronomy, where periodic signals reveal
orbital and rotational structure; and in medical imaging, where
frequency-domain methods reconstruct images from raw scanner data.
Whenever a phenomenon oscillates, Fourier methods provide a powerful
common language for analyzing it.

\section{A Historical Perspective}

Joseph Fourier introduced these ideas while studying the flow of heat
through matter. His original goal was not music, communication, or image
processing; he wanted to understand temperature. The remarkable
discovery was that the same mathematical decomposition appeared in many
unrelated contexts far removed from heat, and what began as a theory of
heat became one of the central tools of modern science.

\section{Looking Back}

This book began with ratios arising from similarity. From similarity
came the trigonometric functions; from trigonometric functions came
rotations; from rotations came amplitwists; from amplitwists came
complex numbers; from complex numbers came derivatives and geometric
series; from geometric series came superpositions of oscillations; and
from superposition came the Fourier idea developed in this chapter. The
path has been long, but it has followed a single theme throughout:
mathematics often progresses by identifying transformations, discovering
what remains invariant under those transformations, and then giving
those invariant structures names. The concepts introduced in this book
are not isolated topics gathered under one title. They are parts of a
single connected framework for understanding pattern, change, and
structure.

\section{Exercises}

\subsection*{Frequency and Harmonics}
\begin{exercise}
Identify the frequencies present in $\sin x + \sin(4x)$.
\end{exercise}
\begin{exercise}
Identify the harmonics appearing in $2\cos x + \cos(5x)$.
\end{exercise}
\begin{exercise}
Which function oscillates more rapidly, $\sin(2x)$ or $\sin(7x)$?
\end{exercise}

\subsection*{Frequency and Period}
\begin{exercise}
Determine the frequency and period of $\sin(4x)$.
\end{exercise}
\begin{exercise}
Determine the frequency and period of $\cos(7x)$.
\end{exercise}
\begin{exercise}
Which function has the shorter period, $\sin(3x)$ or $\sin(5x)$?
\end{exercise}
\begin{exercise}
Find a function whose frequency is twice that of $\cos x$.
\end{exercise}
\begin{exercise}
Find a function whose frequency is five times that of $\sin x$.
\end{exercise}

\subsection*{Superposition}
\begin{exercise}
Sketch $\sin x + \tfrac12\sin(2x)$.
\end{exercise}
\begin{exercise}
Sketch $\cos x + \tfrac13\cos(3x)$.
\end{exercise}
\begin{exercise}
Explain how superposition differs from multiplication of functions.
\end{exercise}

\subsection*{Harmonics and Superposition}
\begin{exercise}
Identify the harmonics present in $f(x) = \sin x + \tfrac12\sin(2x) +
\tfrac14\sin(3x)$.
\end{exercise}
\begin{exercise}
Identify the harmonics present in $f(x) = 2\cos x - \cos(4x)$.
\end{exercise}
\begin{exercise}
Which harmonic has the largest amplitude in $f(x) = 3\sin x + \sin(2x) +
\tfrac12\sin(5x)$?
\end{exercise}
\begin{exercise}
Explain why $\sin x + \sin(2x)$ is more complicated than either term
alone.
\end{exercise}

\subsection*{Complex Exponentials}
\begin{exercise}
Use Euler's formula to expand $e^{i2x}$.
\end{exercise}
\begin{exercise}
Express $e^{i3x}$ in terms of sine and cosine.
\end{exercise}
\begin{exercise}
Explain why $e^{ix} + e^{-ix}$ is real-valued.
\end{exercise}
\begin{exercise}
Show that $e^{-ix} = \cos x - i\sin x$.
\end{exercise}
\begin{exercise}
Show that $e^{ix} + e^{-ix} = 2\cos x$.
\end{exercise}
\begin{exercise}
Show that $e^{ix} - e^{-ix} = 2i\sin x$.
\end{exercise}
\begin{exercise}
Explain why $e^{ix}$ may be interpreted as a rotating unit vector.
\end{exercise}

\subsection*{Finite Fourier Sums}
\begin{exercise}
Express $\cos x + \sin x$ as a finite Fourier sum.
\end{exercise}
\begin{exercise}
Express $2+\cos x+\sin(2x)$ as a finite Fourier sum.
\end{exercise}
\begin{exercise}
Identify the constant term and harmonic terms in $3+2\cos x-\sin(3x)$.
\end{exercise}
\begin{exercise}
How many distinct frequencies appear in $\sin x+\sin(2x)+\sin(4x)$?
\end{exercise}
\begin{exercise}
Explain why every finite Fourier sum is periodic.
\end{exercise}

\subsection*{Orthogonality and Projection}
\begin{exercise}
In Chapter~\ref{ch:ch13-vectors-dot-products}, orthogonal vectors
represented independent directions. Explain the analogous role played
by harmonics.
\end{exercise}
\begin{exercise}
Why is it useful to think of $\cos x$ and $\cos(5x)$ as distinct
directions?
\end{exercise}
\begin{exercise}
Explain how Fourier coefficients resemble vector coordinates.
\end{exercise}
\begin{exercise}
Explain how Fourier decomposition resembles projection onto coordinate
axes.
\end{exercise}
\begin{exercise}
Compare the decomposition of a vector into components with the
decomposition of a signal into harmonics.
\end{exercise}

\subsection*{Applications}
\begin{exercise}
A musical tone contains frequencies $220$ Hz, $440$ Hz, and $660$ Hz.
Which frequencies are harmonics of the fundamental?
\end{exercise}
\begin{exercise}
Why does a violin sound different from a flute even when both play the
same note?
\end{exercise}
\begin{exercise}
Explain how noise-cancelling headphones use superposition.
\end{exercise}
\begin{exercise}
Why might an engineer wish to remove certain frequencies from a signal?
\end{exercise}
\begin{exercise}
Describe one application of Fourier analysis in astronomy.
\end{exercise}
\begin{exercise}
Describe one application of Fourier analysis in medical imaging.
\end{exercise}

\subsection*{Connections Across the Book}
\begin{exercise}
Explain how Chapter~\ref{ch:ch08-rotation-geometry} contributes to the
Fourier idea.
\end{exercise}
\begin{exercise}
Explain how amplitwists lead naturally to complex multiplication.
\end{exercise}
\begin{exercise}
Explain why $e^{i\theta}$ can be viewed simultaneously as a complex
number and a rotation.
\end{exercise}
\begin{exercise}
Explain how the geometric-series formula of
Chapter~\ref{ch:ch18-sequences-series} supports the study of complex
exponentials.
\end{exercise}
\begin{exercise}
Trace the chain rotation $\to$ amplitwist $\to$ complex multiplication
$\to$ Fourier analysis, explaining the role played by each step.
\end{exercise}
\begin{exercise}
Explain why Fourier analysis may be viewed as a study of superposed
rotations.
\end{exercise}

\subsection*{Capstone Problems}
\begin{exercise}
Trace the conceptual chain similarity $\to$ trigonometry $\to$ rotation
$\to$ complex numbers $\to$ Fourier analysis, and explain the role
played by each step.
\end{exercise}
\begin{exercise}
Explain why Fourier analysis can be viewed as a study of superposition.
\end{exercise}
\begin{exercise}
Describe how Chapters~\ref{ch:ch15-amplitwists},
\ref{ch:ch16-complex-euler}, and \ref{ch:ch18-sequences-series}
contribute to the Fourier idea.
\end{exercise}

\subsection*{Challenge Problems}
\begin{exercise}
Use Euler's formula to show that $\cos x = \dfrac{e^{ix}+e^{-ix}}{2}$.
\end{exercise}
\begin{exercise}
Use Euler's formula to show that $\sin x = \dfrac{e^{ix}-e^{-ix}}{2i}$.
\end{exercise}
\begin{exercise}
Explain why every finite Fourier sum can be written using complex
exponentials.
\end{exercise}
\begin{exercise}
Explain why every finite Fourier sum can also be written using only
sines and cosines.
\end{exercise}
\begin{exercise}
A signal contains only frequencies $1$, $3$, and $5$. Write a possible
finite Fourier sum representing such a signal.
\end{exercise}
\begin{exercise}
Explain why the Fourier idea is fundamentally a statement about
representation rather than computation.
\end{exercise}
\begin{exercise}
The chapter began with the Metonic cycle. Explain how the interaction of
two periodic processes can produce a larger apparent cycle.
\end{exercise}
\begin{exercise}
Write a one-page essay describing how the ideas of this book progress
from similarity to Fourier decomposition.
\end{exercise}
